% Optional addition after Proposition prop:period-input-audit.
% Independent audit: 2026-09-05. No canonical edits made.
\begin{proposition}[exact admissible period constants]
\label{prop:period-exact-admissible-constants}
Let \(\tau,\mu,\beta^{\mathrm{part}}\) be the cusp-regular functions
constructed in \cref{prop:period-input-audit}, with all its marked
coordinates and transformation laws retained.  The descended function
\[
 D_{\mathrm{part}}=
 \operatorname{Im}\beta^{\mathrm{part}}
 -\frac{6(\operatorname{Im}\mu)^2}{\operatorname{Im}\tau}
\]
has an attained finite maximum \(M\) on
\(B\setminus\{p_0\}\).  The exact set of constants for which
\(\beta=\beta^{\mathrm{part}}+c_0\) makes the marked real period map an
isomorphism at every point is
\[
 \{c_0\in\mathbb C:\operatorname{Im}c_0<-M\}.
\]
\end{proposition}

\begin{proof}
The preceding proposition proves continuity on
\(B\setminus\{p_0\}\), including its elliptic points, and convergence
\(D_{\mathrm{part}}\to-\infty\) at \(p_0\).  Choose a base point
\(b_*\).  Outside a compact set containing \(b_*\), the function is
smaller than \(D_{\mathrm{part}}(b_*)-1\).  Its maximum on that compact
set is therefore its maximum on the entire base.  A constant shift gives
\(D=D_{\mathrm{part}}+\operatorname{Im}c_0\), so \(D<0\) everywhere
exactly for \(\operatorname{Im}c_0<-M\).  At equality, \(D=0\) at a
maximum point.  For \(\operatorname{Im}c_0>-M\), that point has positive
\(D\), whereas the cusp still has negative \(D\).  A path toward the
cusp contains a zero.  Since the exact ordered real determinant is
\(\operatorname{Im}\tau\,D\), these and only these excluded constants
make the period matrix singular somewhere.
\end{proof}

\begin{remark}[the cusp extension is essential]
For the same \(\tau,\mu,\beta\), the function
\(\widetilde\beta=\beta+i\,t\circ\pi\) is holomorphic on
\(\mathfrak h_z\) and satisfies all the same affine transformation laws,
because \(t\circ\pi\) is invariant.  However,
\(\widetilde D=D+\operatorname{Re}t\).  On the actual ray \(t=R\to+\infty\),
\(j(\tau)=1728R\) gives
\[
 q=(1728R)^{-1}(1+O(R^{-1})),\qquad
 D(R)=-\frac{\log(1728R)}{2\pi}
       +\operatorname{Im}b(0)+o(1),
 \quad b=\beta+\tau.
\]
Consequently \(\widetilde D(R)\to+\infty\), and no additive constant
makes \(\widetilde D\) negative everywhere.  This function is excluded
by the actual construction because
\(\widetilde\beta+\tau=b+i t\) has a cusp pole.  Thus holomorphic
solvability on the open base alone does not supply the required uniform
bound; the proved cusp extension supplies it.
\end{remark}
